\section{Scope, Mediation, and Limitations}
\label{sec:discussion}

\subsection{Perimeter}
\label{sec:scope}

What we tested is the observable on-chain channel. A reflexive
leverage-cycle spiral, in the strict sense, plausibly lives on the
centralized perpetual-futures venue, a compartment we do not observe and
therefore neither test nor refute. Our null is not a denial of
centralized-exchange cascades. It is a result about a different object, three
to four orders of magnitude smaller. The scale gap carries this reading. The
typical daily DeFi flow is four to six orders of magnitude below the largest
single-day centralized cascade on record (roughly nineteen billion dollars,
October 2025; Section~\ref{sec:size}), and the largest crashes of the period (Terra, FTX,
October 2025) are uniformly attributed to centralized, perpetual, or systemic
dynamics, never to on-chain DeFi. We are deliberate about the evidentiary
status of that attribution. The weight of evidence favors the centralized
venue, but the causal claim is associative on both sides: event studies,
connectedness measures, and practitioner magnitudes, the very class of
non-identified evidence this paper otherwise declines. The formulation we
prefer is that the centralized venue is the most plausible locus, not an
established cause, and that our DeFi design cannot adjudicate it. The
supporting magnitudes are labeled as event-study and practitioner figures
throughout, and the measured on-chain price impact remains minuscule and temporary \citep{tian2025}, real locally but too small to aggregate.

\subsection{The null in the crypto leverage-effect debate}
\label{sec:mediation}

The field is split. In spot data at lower frequencies the leverage effect
is inverted or absent \citep{baur2018,aharon2023}, while the cascade
narrative, in which liquidations drive the downside, lives in a different
venue, frequency, and jump regime. Our null is the tail-response counterpart
of the inverted-or-absent spot effect, transposed to a new object: net of
volatility, the DeFi shock predicts large up-moves and large down-moves
equally. It is not in tension with the cascade results, which inhabit an
intraday, jump-driven, centralized regime that an hourly on-chain shock does
not capture. The reconciling axis is the decomposition of \citet{huang2022},
in which the diffusive component behaves conventionally while jumps invert, so
the net sign flips with the mix. That is why spot and cascade findings coexist
without contradiction. Our contribution to the debate is one clean empirical
point that sits squarely in the both-tails-respond-more-than-the-mean camp
\citep{ando2022}. We do not adjudicate the true sign of the crypto leverage
effect; we place our null within it. One apparent counter deserves a direct
word. Cross-coin tail dependence is, if anything, stronger on the
\emph{upside} \citep{nguyen2020}, but that is co-dependence across coins, not
a single asset's response to a shock, and it is consistent with our finding of
no downside specificity.

\subsection{Scope of the null}
\label{sec:caution}

The null must be stated precisely. At the observable on-chain scale and over
this period, there is no economically significant downside amplification. That
is not the same as saying there is none, or that there will be none. The
conditional is explicit: \emph{if} the DeFi share of crypto leverage were to
grow materially, the channel \emph{could} become aggregate-relevant. The
weight of that statement is carried by three legs, size
(Section~\ref{sec:size}), mechanism
(Sections~\ref{sec:clearing}--\ref{sec:oracle}), and the official-sector
consensus that DeFi is not yet systemic but bears watching as it grows
\citep{fsb2023,aramonte2021,imf2021,ecb2022}, never by the whisper. The
conditionality is not rhetorical: on-chain leverage \citep{heimbach2024} and
cross-protocol interconnection \citep{badev2023} are each time-varying, and
the ecosystem's measured fragility moves with them \citep{bertomeu2024}. This is why we bound the null
rather than assert a zero. An equivalence bound
(Section~\ref{sec:boundednull}) excludes only a downside asymmetry larger than
roughly \MdeEightyExc{} per unit of shock, which is what makes the null
respectable and falsifiable rather than a bare failure to reject
\citep{altman1995,ioannidis2017}.

The one-percent whisper bounded in Section~\ref{sec:boundednull} is, in
addition, governed at $\tau = 0.01$ by a non-Gaussian limiting distribution
that makes standard inference unreliable to begin with
\citep{chernozhukov2005,chernozhukov2011}. That is one more reason it does not
carry the conclusion, which rests on size, mechanism, and consensus.

\subsection{Limitations}
\label{sec:limitations}

We list the limitations, say why each is one, and state what was done about
it. \emph{(1) Hourly aggregation.} The shock is aggregated to the hour, so
within-hour sequencing is lost and the impact hour mixes trigger and
consequence; we adopt a predictive register and read across horizons, but
no tick-level sequencing is available. \emph{(2) BTC--ETH co-movement.} The
impact response replicates on Bitcoin (Section~\ref{sec:btc}:
$\beta^{\text{BTC}}_0(0.01) = \BtcBetaQoneHzero$ against
$\BetaQoneHzero{}$ for ETH), so the shock is partly a common-crypto stress
proxy rather than purely ETH-specific; the cross-asset placebo is precisely
the diagnostic for this, and the result is what motivates retiring the
specificity claim rather than defending it. \emph{(3) Short-horizon reverse
causality.} \PLiqWorstOne\% of the worst-1\% hours contain contemporaneous
liquidations against \PLiqBestOne\% of the best, so a crash mechanically
triggers liquidations and the short-horizon relation is bidirectional. Hence
the non-causal register, with the one-hour lag and horizon reading as
mitigation. \emph{(4) No centralized exchange.} The perimeter is on-chain
only, so the plausible locus of the spiral is structurally out of reach.
\emph{(5) Mean, not quantile, residualization.} The orthogonalized shock is
the residual of an OLS, not a quantile, regression. We acknowledge this as a
refinement, with the raw in-regression controls serving as the canonical
channel and the orthogonalized shock reserved for the placebo.
\emph{(6) Estimated denominators.} The size ratio rests on
centralized-reported, partly wash-inflated turnover figures, so it is an
order-of-magnitude bound; the numerator is the only hard figure, and the
inequality survives one to two orders of rescaling. \emph{(7) No structural
identification.} Every estimate is a predictive association; converting them
to causal statements would require exogenous variation in liquidation
timing, the most immediate direction for future work. Limitations (1)--(3)
bound causality, (4) bounds scope, and (5)--(7) bound precision; none
invalidates the beam, and together they fix the perimeter of the claim.